\chapter{Snow basics}

\subsection{Direct and indirect methods for snow measurement}
Knowledge of snow cover properties is essential for assessing solid water storage in the snowpack, a critical factor influencing agriculture and hydroelectric power generation, as well as for avalanche risk evaluation. 

Historically, snow metrics were recorded manually, requiring an operator to conduct manual measurements on a daily basis whenever feasible. 
To ensure continuous data collection even in remote areas, automated weather stations (AWS) have been introduced, and are steadily replacing manual ones as well. An automated monitoring network provides a comprehensive overview for hazard assessment and for studying snow dynamics over decades. 
However, despite the significant number of AWS currently in operation, achieving high-resolution coverage remains challenging due to high costs and maintenance issues. Consequently, certain data can only be acquired via remote sensing, by means of Earth-observing satellites.

While remote sensing provides spatially distributed data across entire regions, manual and automated measures are taken at specific locations, which are chosen according to strict criteria.
To ensure the reliability of these in situ observations, the monitoring site must match the features of the surrounding terrain.
In alpine regions, monitoring sites must be clear of large boulders and protected from excessive sun exposure, which can cause premature snowmelt and yield unrepresentative data.
It is also important to avoid locations prone to preferential snow accumulation driven by local wind patterns.
In forested areas, instruments should be placed within openings to ensure that snowfall reaches the ground and accumulation can take place unbiased.

These guidelines represent only the baseline requirements for reliable snow measures. A more comprehensive analysis of the operational details is provided in the following sections.


\subsubsection{Manual measures}

Manual snow depth (HS) records are an effective tool for assessing long-term trends, as they can span over several decades. Their longevity is due to a straightforward measuring procedure, requiring only a stake or an extendible graduated rod. 

The snow stake should be vertical, painted white to minimize melt and fixed to the ground before accumulation begins. It is required to make snow depth observations from a sufficient distance to avoid any deformation of the snowpack.
Because the stake acts as an obstacle to wind and snow redistribution, localized mounds can form around it. Under such circumstances, operator must visually estimate the undisturbed snow surface to recover the actual HS value.

For a non-fixed observation, the snow depth value should be evaluated at the same location each time. The observer needs to ensure that the device, manually inserted in the snowpack, reaches the base surface without penetrating any soil thus overestimating the actual HS. The rod has to be vertical, and an estimate of the slope angle should be paired to every measure.
% Immagine per snow stake e snow rod
\newline
\newline
Closely linked to snow depth is snowfall depth (HN), which can only be obtained manually. To avoid wind-induced bias this measurement should be carried out in a sheltered location, so that snow accumulates undisturbed on an artificial surface for the prescribed interval (usually 24 hours). Just before clearing the surface, thereby starting the next observation period, a ruler is inserted vertically into the freshly formed layer to obtain a depth measurement.
\newline
\newline
Although HS and HN define the vertical extension of the snowpack, hydrological applications rely primarily on the snow water equivalent (SWE). As shown in (cite equation), SWE is directly linked to snow depth through the bulk density of the snowpack. Although various manual measurement techniques exist, this section focuses on the two primary methods, snow pits and snow tubes.

A snow pit consists of a manually-dug hole up to a reference surface, usually the ground. Initially, the observer performs a stratigraphic analysis to deepen his knowledge of the internal structure of the snowpack. The actual SWE measurement is conducted by inserting a sampler of known volume perpendicularly into the exposed snow face; the sample is then weighted to assess snow density. 
This procedure is usually repeated at multiple depths, either for each identified layer or at fixed intervals, such as every 20 centimetres.
This method enables the observer to reconstruct the snow density profile which, when combined with depth information, yields both the layer-by-layer and total snow water equivalent at that specific location. Clearly, the snow pit technique is disruptive and needs to be performed at a different site each time.
% Immagine schematica per come debba essere effettuata una campagna con snow tube

If the measurement campaign is focused on a spatial survey rather than a detailed, locally defined inspection, the logical choice is utilizing a snow tube, because it allows for rapid sampling of the entire snow column to directly retrieve the total SWE, sacrificing layer-specific information.
The snow tube is inserted vertically in the snowpack until it reaches the ground. An estimate of the extracted core volume is made, and the average density is then computed by weighting the snow column. The conversion to SWE is then trivial.
This device should be used where the snow water equivalent is less spatially homogeneous, such as locations with heterogeneous exposure to both wind patterns and solar radiation.



\subsubsection{Automatic measures}

Automatic weather station (AWS) networks do not require manual intervention, except for routine maintenance and instrument calibration. An AWS provides continuous observations even in remote locations where human operation would be unfeasible.
Typically, an alpine automatic weather station is equipped with a variety of sensors to provide a highly comprehensive overview of environmental conditions. 

Automated HS measurements can be conducted using instruments based on either sonic or optical (laser) technology.
Sonic sensors transmit an ultrasonic pulse towards the snow surface and detect the returning echo. Since the speed of sound is temperature-dependent, a temperature reading is also required to ensure accuracy. On the other hand, laser instruments determine the target surface distance by analyzing signal interference. These measurements are actually much more precise, but they come at a higher energy cost.
In both scenarios, the sensor measures the distance to the upper layer of the snowpack, from which the snow depth is then retrieved 

\begin{equation}
    \mathrm{HS} = \mathrm{Zero\ depth\ distance} - \mathrm{Distance\ to\ target}
    \label{eq:auto_HS}
\end{equation}

The zero depth distance must be determined before snow accumulation begins, and is a key part of instrument calibration. 
To obtain reliable measurements, the HS sensor should be installed at a sufficient height above the maximum anticipated snow depth to prevent burial during winter, which would result in data loss. In addition, minimizing interference with the support structure is essential, as the mast acts as a wind obstacle and can induce measurement artifacts. Consequently, HS sensors are typically installed on a horizontal arm that stretches outward, ensuring that the target area remains at a sufficient distance from the mounting pole.

One limitation of HS series is that, from a hydrological perspective, they do not directly quantify the amount of water stored within the snowpack. Several techniques have been developed for automated SWE monitoring, the most common of which based on snow pillows. 
This instrument consists of a rubber bladder filled with an antifreeze liquid mixture  and coupled with a pressure transducer, which provides the SWE estimate. 
The snow pillow is installed flush with the ground surface, thereby supporting the weight of the entire snowpack. As snow accumulates on top of the pillow, the hydrostatic pressure increases and this variation is transduced into snow water equivalent values. 
Snow pillows were first introduced in the late 1960s, and have proved to be an useful tool for SWE measurements. 
% cita studio americano monte hood 1965
However, they are currently being replaced by snow scales due to the toxicity of the antifreeze fluid, which poses an environmental hazard to the installation site.